\documentclass[11pt]{article}
\usepackage[margin=1in]{geometry}
\usepackage{amsmath,amssymb}
\usepackage[hidelinks]{hyperref}

\title{Disproof of Conjecture 00000001024\\
\large When the prescribed BCH distance is attained infinitely often}
\author{}
\date{}

\begin{document}
\maketitle

Let \(N=2^m-1\) and let \(\alpha\) generate
\(\mathbb F_{2^m}^{\times}\). In the primitive narrow-sense binary
BCH convention \cite{brc}, prescribing zeros
\(\{1,\ldots,\delta\}\) means imposing
\[
  C_{m,\delta}
  =\{c(X)\in\mathbb F_2[X]:\deg c<N,\quad
                       c(\alpha^k)=0\quad(1\le k\le\delta)\}.
\]
The actual root set includes the Frobenius closure of these prescribed
zeros. Write \(d=\delta+1\) for the proposed minimum distance. Since
this is a linear code, its minimum distance is the minimum nonzero
Hamming weight.

\section*{Odd \(\delta\) always forces a larger distance}

When \(\delta\) is odd, \(d\) is even and \(1\le d/2\le\delta\).
For every word \(c\in C_{m,\delta}\), its binary coefficients give
\[
  c(\alpha^d)=c(\alpha^{d/2})^2=0.
\]
Thus \(1,\ldots,d\) are root exponents, including the next exponent
after the prescribed interval.

To see the distance implication, suppose a nonzero word of weight
\(t\) vanishes at each of \(\alpha,\ldots,\alpha^t\). Its support
positions \(j_1,\ldots,j_t\) are distinct in \(\{0,\ldots,N-1\}\).
With \(x_i=\alpha^{j_i}\), the syndrome equations are
\(\sum_{i=1}^t x_i^k=0\) for \(k=1,\ldots,t\). This matrix is
Vandermonde after factoring \(x_i\) from column \(i\); its determinant
is nonzero because the \(x_i\) are distinct and nonzero. It cannot
annihilate the vector of all ones. This is the BCH lower bound for
consecutive zeros starting at \(1\).

For each fixed odd \(\delta\ge1\) and every sufficiently large \(m\),
applying this argument to the forced roots excludes every nonzero
word of weight at most \(d\). The code has a nonzero word for such
\(m\), since the least common multiple of the minimal polynomials
of the prescribed roots has degree at most \(m\delta<N\).
In particular the distance cannot be
\(\delta+1\) for infinitely many \(m\).

For a concrete family, take \(\delta=1\). A codeword has zero sum of
its locators. A single nonzero locator cannot sum to zero, and two
distinct locators cannot sum to zero in characteristic two. For
every \(m\ge2\), take distinct nonzero \(a,b\) in the field. The
three locators \(a,b,a+b\) are distinct and nonzero, and they sum
to zero. Each is a power of \(\alpha\), so they give a word of
weight three. The exact distance is therefore three for every
\(m\ge2\), whereas the conjecture gives two.

\section*{Even \(\delta\) does attain equality infinitely often}

For even \(\delta\ge2\), \(d=\delta+1\) is odd. Since \(2\) is a unit
modulo \(d\), the relation \(d\mid 2^m-1\) holds for all positive
multiples of the order of \(2\) modulo \(d\). For any sufficiently
large such \(m\), take
\[
  c(X)=\sum_{j=0}^{d-1} X^{jN/d}.
\]
It has weight \(d\), with all exponents distinct and less than \(N\).
The element \(\zeta=\alpha^{N/d}\) has order \(d\). For
\(1\le k\le d-1\), the ratio \(\zeta^k\) is different from one
but has \(d\)-th power one, so the geometric sum gives
\(c(\alpha^k)=\sum_{j=0}^{d-1}(\zeta^k)^j=0\).
This word belongs to \(C_{m,\delta}\), and the Vandermonde argument
for its prescribed roots gives distance at least \(d\). Thus the
distance is \(d\) for all these \(m\).

\section*{Lean verification}

The Lean~4.19.0 project models binary support words by lists of
distinct locators. It proves the exact distance-three result for
every \(m\ge2\) from the characteristic-two additive laws and a
complete enumeration of \(2^m-1\) nonzero locators. It also checks the
Frobenius doubling identity on arbitrary supports and shows that
every odd cutoff forces the next syndrome to vanish. The primitive
field interpretation, Vandermonde step, and odd-design-distance
construction are proved above. Run \verb|lake build| in \verb|lean/|
to verify the formal results.

\begin{thebibliography}{1}
\bibitem{brc}
R.~C. Bose and D.~K. Ray-Chaudhuri,
\emph{On a class of error correcting binary group codes},
Information and Control \textbf{3} (1960), 68--79.
\href{https://doi.org/10.1016/S0019-9958(60)90287-4}{DOI}.
\end{thebibliography}

\end{document}
